\documentclass[a4paper,12pt]{article}
\input{LaTeX-ToolKit/template.tex}
\begin{document}
\title{Band Theory}
\author{沈威宇}
\date{\temtoday}
\titletocdoc
\sct{Band Theory (能帶理論)}
\ssc{Basics}
\sssc{Electronic band structure or band structure}
The electronic band structure (or simply band structure) of a solid describes the range of energy levels that electrons may or may not have within it.
\sssc{(Energy) band (能帶)}
Atomic orbitals from multiple atoms combine to form molecular orbitals, such as bonding and anti-bonding orbitals. When a vast number of atoms—typically on the order of $10^{20}$ or more—come together to form a solid, the energy differences between these molecular orbitals become extremely small. Consequently, a cluster of closely spaced energy levels can be treated as a continuous energy band.
\sssc{Fermi level (費米能階)}
The Fermi level of a solid-state body is the thermodynamic work required to add one electron to the body. For Cu, $E_F=7.0$ eV.
\sssc{Work function (功函數)}
For a material with Fermi level $E_F$ and energy of an electron infinitely far in vacuum $E_{\tx{vac}}$, the work function is
\[W=E_{\text{vac}}-E_F.\]
Work function is related with the threshold frequency $f_0$ of photoelectric effect by
\[W=hf_0.\]
\ssc{Metals}
\sssc{Bands}
A partially filled band always exist, in which electrons can wander around and conduct current, called free electrons.
\sssc{Number density of electrons in partially filled band}
The number density (number per unit volume) of electrons in partially filled band $n$ is the material's mass density multiplied by the material's number of valence electrons per atom multiplied by Avogadro's number divided by the material's molar mass.
\sssc{Number density of states of the allowed energy level per unit volume per unit energy interval}
The number density of states of the allowed energy level per unit volume per unit energy interval $N(E)$ is
\[N(E)=\frac{8\sqrt{2}\pi m_e^{\pht{e}3/2}}{h^3}\sqrt{E}.\]
\sssc{Occupancy probability}
The occupancy probability $P(E)$ is the probability that a given available state will be occupied by an electron
\[P(E)=\frac{1}{e^{\frac{E-E_F}{kT}}+1}\]
\sssc{Number density of occupied states per unit volume per unit energy interval}
The number density of occupied states per unit volume per unit energy interval $N_0(E)$ is
\[N_0(E)=N(E)P(E).\]
\sssc{Finding Fermi energy}
$E_F$ can be found by solving
\[n=\int_0^{E_F}N_0(E)\dd{E}.\]
The result is
\[E_F=\qty(\frac{3}{16\sqrt{2}\pi})^{2/3}\frac{h^2}{m}n^{2/3}\approx\frac{0.121h^2}{m}n^{2/3}.\]
\sssc{Thermal effect}
The change of Fermi level on temperature change is small and often negligible for metals.
\ssc{Semiconductors and Insulators}
\sssc{Bands}
\bit
\item\tb{Valence band (價帶)}: The lower-energy, filled band, in which electrons are constrained in covalent bonds and can't wander away freely or conduct current.
\item\tb{Conduction band (導帶)}: The higher-energy band that is empty or partially empty, in which electrons, if exist, can wander around and conduct current, called free electrons.
\item\tb{Band gap (能隙)}: The gap between conduction band minimum and valence band maximum. The energy difference of the gap is called band-gap energy $E_g$. For silicon, $E_g=1.12$ eV.
\eit
\sssc{Insulator}
A material with empty conduction band and high energy band such that virtually no transition of electron from valence band to conduction band occurs.
\sssc{Semiconductor}
A material with lower energy band such that transition of electron from valence band to conduction band can occur under room temperature, that is, whose conductivity lies between that of conductors and inductors.

There are two kinds of semiconductors: single-element semiconductors, such as silicon (IV) and germanium (IV), and compound semiconductors, such as gallium arsenide \ce{GaAs} (III-V), aluminium gallium arsenide \ce{AlGaAs} (III-V), indium phosphide \ce{InP} (III-V), Aluminium gallium indium phosphide \ce{AlInGaP} (III-V), zinc oxide \ce{ZnO2} (II-VI), and titanium doxide \ce{TiO2} (II-VI). Compound semiconductors are useful in special electronic circuit applications as well as in applications that involve light, such as light-emitting diodes (LEDs). Germanium was used in the fabrication of very early transistors (late 1940s, early 1950s). It was quickly supplanted by silicon, on which today’s integrated-circuit technology is almost entirely based. For this reason, we will deal mostly with silicon devices.
\ssc{Effective mass, holes, and carriers of semiconductors}
\sssc{Energy and effective mass (有效質量) of electrons}
The energy of an electron in a semiconductor is not the simple square relation of that of a free electron and wave number but a complicated periodic function. However, near an edge of a band, it can be approximated with:
\[E(k) \approx E_0 + \frac{\hbar^2 k^2}{2m^*}\]
where $E_0$ is the energy of the edge, $k$ is wave number, and $m^8$ is a constant called effective mass given by
\[m^*=\hbar^2\qty(\frac{\mathrm{d}^2E(k))}{\mathrm{d}k^2})^{-1}\]
where $\dv[2]{E(k)}{k}$ is taken at the edge. Effective mass is usually written as multiple of mass of electron.
\sssc{Effective mass of electrons at difference levels}
\bit
\item At conduction band minimum, $\dv[2]{E(k)}{k}>0$ and $m^*>0$. These electrons behave like free electrons but with mass $m^*$.
\item At valence band maximum, $\dv[2]{E(k)}{k}<0$ and $m^*<0$. These electrons move to the opposite direction of a force acting on it.
\eit
\sssc{(Electron) hole (（電子）電洞/空穴)}
A hole $h^+$ is an absence of an electron in valence band.

Let the effective mass of electrons at valence band maximum be $-m^*$ and the concentration of holes be $p$. Consider a nearly-filled valence band, the current of it equals to the current of a filled valence band, which is always zero, subtracting the current that would exist due to the electrons that would fill the holes if there are not the holes, which is in the opposite direction of the force. Thus, the current is in the direction of the force and equals to that due to particles with charge $e$ and mass $m^*$ at concentration $p$.

Thus, we define hole as a quasiparticle with charge $e$ and effective mass $m^*$, which is positive.
\sssc{Carriers (載子) of semiconductors}
Thus, there are two types of carriers in semiconductors mathematically:
\bit
\item free electron at conduction band minimum with charge $-e$ and a positive effective mass $m_e^*$, and
\item holes at valence band maximum with charge $e$ and a positive effective mass $m_h^*$.
\eit
The concentration of free electrons is denoted $n$; the concentration of holes is denoted $p$.
\sssc{Carrier generation}
When an electron transits from valence band to conduction band, a free electron in conduction band and a hole in valence band are generated, and an energy of band-gap energy $E_g$ is consumed.

Since they are generated in pairs, an intrinsic (undoped) semiconductor has equal concentration of free electrons and holes.
\sssc{Carrier recombination (載子複合) and electroluminescence (電致發光/電激發光)}
When an electron transits from conduction band to valence band, a free electron in conduction band and a hole in valence band are cancelled, and an energy of band-gap energy $E_g$ is released in the form of a photon with frequency $f_h=\frac{E_g}{h}$, called electroluminescence.

Carrier recombination occurs naturally and at a rate proportional to the number of free electrons and holes.

Semiconductors with different band-gap energy can release light of different frequency. This effect is used in light emitting diode (LED) (發光二極體).
\sssc{Exciton (激子)}
A quasiparticle with no charge of a pair of free electron and hole in bounded state due to electrostatic force.
\sssc{Photovoltaic effect (光生伏打效應/光伏效應)/Internal photoelectric effect (內光電效應)}
Carrier generation caused by light.

Similar to the photoelectric effect, if the frequency of the incident light is greater than or equal to the material's threshold frequency $f_h=\frac{E_g}{h}$, electron-hole pairs are generated immediately; if the frequency of the incident light is less than this threshold, such pairs are never produced.

This effect is used in solar cell, photoresistor, and light-controlled relay.
\sssc{Thermal generation and carrier concentration at thermal equilibrium}
Carrier generation due to thermal energy. The higher the temperature, the more thermal generations occur.

Thermal equilibrium is the state where thermal generation rate equals to carrier recombination rate, at which
\[np=n_i^2\]
where $n_i$ denotes the number of free electrons and holes in a unit volume at a given temperature given by
\[n_i=BT^{\frac{3}{2}}e^{-\frac{E_g}{2kT}}\]
where $B$ is a material-dependent parameter that is $7.3\times 10^{15}$ cm$^{-3}$K$^{-\frac{3}{2}}$ for silicon, $T$ is temperature in K, and $k$ is Boltzmann's constant $8.62\times 10^{-5}$ eV/K.

For intrinsic semiconductor, that is,
\[n=p=n_i\]
At 0 K, no thermal generation occurs, so an intrinsic semiconductor has no carrier and behaves as an inductor at thermal equilibrium.
\ssc{Doping}
\sssc{Doping (摻雜)}
Doping is the intentional introduction of impurities into an intrinsic/undoped/pure semiconductor in order to change its properties. The doped material is referred to as an extrinsic/doped semiconductor.
\sssc{n-typed (n 型) semiconductors}
Semiconductor with doping the 
摻雜易於失去電子的元素，稱donor，引入更多自由電子。對於類金屬半導體通常是在高純度半導體中加入少量第15族元素，如 P、As。費米能階上移而靠近導帶，導電性增加。電子是major carriers，電洞是minor carriers。
\sssc{p 型（p-typed）半導體}
摻雜易於獲得電子的元素，稱acceptor，引入更多電洞。對於類金屬半導體通常是在高純度半導體中加入少量第13族元素，如 B、Al。費米能階下而靠近價帶，導電性增加。電子是minor carriers，電洞是major carriers。
\ssc{Intrinsic semiconductors}
\ssc{p–n 接面（p–n junction）}
\sssc{Definition}
A p-n junction is a single semiconducting crystal with one end doped to form p-type material and the other end doped to form n-type material. The two types meet at a junction plane. At thermal equilibrium, the following occur at the junction plane:
\bit
\item Majority carriers diffuse across the plane, producing a diffusion current.
\item Minority carriers are swept across the plane, forming a drift current.
\item A depletion zone forms at the plane. It is normally a region of high resistivity, but can be decreased by decreasing its width.
\item A contact potential develops across the depletion zone. The larger it is, the less the diffusion current is. The drift current is not affected by it.
\eit
\sssc{Junction rectifier}
A p-n junction conducts electricity better for one direction of an applied potential difference (forward biased) than for the opposite direction (back biased), and thus the device can serve as a junction rectifier.

In forward-bias connection, the p side becomes more positive and the n side becomes more negative, thus decreasing the height of the potential barrier, increasing the diffusion current, narrowing the depletion zone, and decreasing its resistivity.

In back-bias connection, the p side becomes less positive and the n side becomes less negative, thus increasing the height of the potential barrier, decreasing the diffusion current, widening the depletion zone, and increasing its resistivity.
\sssc{發光二極體（Light emitting diode, LED）}
將順向偏壓通入半導體二極體，使depletion zone極窄，使其中大量發生載子複合而電致發光。The frequency $f$ of the emitted light is
\[E_g=hf.\]

早期研發出的紅色和綠色單色發光二極體用途較受限，1992年赤崎勇、天野浩與中村修二以氮化鎵材料發明藍光 LED，始可用光的三原色調配任意所需的顏色，尤其用作照明的白光，較傳統白熾燈泡與日光燈更省電、體積更小、壽命更長，共同獲2014年諾貝爾物理獎。
\sssc{有機發光二極體（Organic light-emitting diode, OLED）}
有機半導體製成的二極體，原料如「三(8-羥基喹啉)鋁/三[8-羥基(苯[b]氮雜環己-1,3,5-三烯)]鋁（tris(8-hydroxyquinolinato) aluminium(III)）\ce{Al(C9H6NO)3}」或「聚對苯乙炔（Poly(p-phenylene vinylene), PPV） [-c1ccc(cc1)C=C-]$_n$」。
\sssc{Photodiode}
Shining light on a suitably arranged p-n junction can that includes the junction. This is the basis for the photodiode.
\sssc{Junction laser}
In a forward-bias connection connected to a source, there are many electrons in the conduction band of the n-type material and many holes in the valence band of the p-type  material. Thus, there is a population inversion for the electrons; that is, there are more electrons in higher energy levels than in lower energy levels. When a single electron moves from the conduction band to the valence band, it can release its energy as a photon. This photon can stimulate a second electron to fall into the valence band, producing a second photon by stimulated emission. In this way, if the current through the junction is great enough, a chain reaction of stimulated emission events can occur and laser light can be generated. To bring this about, opposite faces of the p-n junction crystal must be flat and parallel, so that light can be reflected back and forth within the crystal. Junction lasers are built into compact disc (CD) players, where, by detecting reflections from the rotating disc, they are used to translate microscopic pits in the disc into sound. They are also much used in optical communication systems based on optical fibers.
\sssc{光伏電池（Photovoltaic cell）/太陽能電池（Solar cell）}
通常由一個或多個 p–n 接面組成，吸收光子並輸出直流電。種類如：
\begin{itemize}
\item 單結晶矽太陽能電池：高效率但成本較高。
\item 多晶矽太陽能電池：效率較低但成本較低。
\item 薄膜太陽能電池：輕薄柔軟，適用於特定場景。
\end{itemize}
效率提升方法：
\begin{itemize}
\item 減少表面反射：透明導電層和抗反射層減少光子的反射。
\item 改進材料：新材料如有機太陽能電池和混合鈣鈦礦太陽能電池提高效率。
\item 多接面結構：提高光吸收和轉換效率。
\item 雷射細工：提高電池表面光陽極的粗糙度，增加光吸收。
\end{itemize}
\sssc{場效電晶體（Field-effect transistor, FET）}
FET 具有四個端子：源極（source）、閘極（gate）和汲極（drain）、體（body）/基（base）/塊體（bulk）/基板（substrate）。通過在閘極施加電壓，可以改變源極與汲極之間的電導率，從而控制電流，分為需加閘極電壓方能導通源極與汲極的增強模式（enhancement-mode），較常見；與預設導通源極與汲極，加閘極電壓可減弱或關閉它的空乏模式/減少模式（depletion-mode）。

下討論 metal-oxide-semiconductor-field-effect transistor (MOSFET)。其僅使用電子與電洞的其中一種作為電荷載子，前者稱 n-通道（n-channel），後者稱 p-通道（p-channel）。源極與汲極下方的摻雜半導體具有該種載子，在增強模式中預設不相連而斷路，但加閘極電壓可以引入載子使相連而導通；在空乏模式中預設相連而導通，但加閘極電壓可以引入相反的載子使變窄增加電阻，繼而不相連而斷路。
\end{document}
